% GENERATED FILE. Edit canonical lessons, then run npm run generate:books.
% canonical-source-hash: fnv1a64:f8f996c56004ec97
% canonical-lessons: IT-C274-nel-frattempo, IT-C274-contemporaneamente, IT-C274-non-appena, IT-C274-a-un-certo-punto, IT-C274-d-un-tratto

\chapter{Nel frattempo}
\label{ch:it-nel-frattempo}
\hlchaptermodality{274}

\begin{chapteropening}
\textbf{By the end of this chapter:} \emph{I can say that two things happened at the same time, or that something broke in all of a sudden.}

Complete the last lesson of chapter 274: I can say that two things happened at the same time, or that something broke in all of a sudden.
\end{chapteropening}

\section[nel frattempo]{nel frattempo — in the meantime, meanwhile}

\label{lesson:IT-C274-nel-frattempo}



\begin{quote}
Before the new one: say the Italian for subsequently, then the Italian for after which.
\end{quote}

\subsection*{You'll want to know: nel frattempo}
\textbf{nel frattempo} — \textquotedblleft{}in the meantime, meanwhile\textquotedblright{}.

\begin{quote}
Lui è andato a prendere il pane; nel frattempo io ho fatto il caffè.
\end{quote}

\begin{quote}
\emph{He went to get the bread; in the meantime I made the coffee.}
\end{quote}

Close to \emph{intanto}, a little more formal.

\begin{cousinweb}
\emph{frattempo} joins \emph{fra} \textquotedblleft{}between\textquotedblright{} and \emph{tempo} \textquotedblleft{}time\textquotedblright{}, which is inherited from Latin \emph{tempus} \textquotedblleft{}time\textquotedblright{} (Lewis \& Short): the time in between.
\end{cousinweb}

\begin{grammarlens}[title={What you've built}]
One more word for telling, explaining and describing at B1.
\end{grammarlens}

\subsection*{Your turn}
\begin{itemize}
\raggedright
  \item \emph{Say it:} \emph{nel frattempo}
  \item \emph{Say it:} \emph{nel frattempo}, once more
  \item \emph{Say it:} a sentence of your own with \emph{nel frattempo}
  \item \emph{Recall:} say the Italian for subsequently, then the Italian for after which, then say \emph{nel frattempo} again
\end{itemize}

\begin{tcolorbox}[breakable,colback=teal!4,colframe=teal!35!black,title={Before you move on}]
What does nel frattempo mean? (\textquotedblleft{}In the meantime, meanwhile\textquotedblright{}.) Say it once more.
\end{tcolorbox}
\section[contemporaneamente]{contemporaneamente — at the same time, simultaneously}

\label{lesson:IT-C274-contemporaneamente}



\begin{quote}
Before the new one: say the Italian for after which, then the Italian for in the meantime.
\end{quote}

\subsection*{You'll want to know: contemporaneamente}
\textbf{contemporaneamente} — \textquotedblleft{}at the same time, simultaneously\textquotedblright{}.

\begin{quote}
Due telefoni hanno suonato contemporaneamente.
\end{quote}

\begin{quote}
\emph{Two phones rang at the same time.}
\end{quote}

Built on \emph{contemporaneo}, ‘of the same time’.

\begin{cousinweb}
\emph{contemporaneamente} is built on \emph{contemporaneo}, a learned borrowing from Latin \emph{contemporāneus} \textquotedblleft{}of the same time\textquotedblright{}, from \emph{con-} \textquotedblleft{}together\textquotedblright{} and \emph{tempus} \textquotedblleft{}time\textquotedblright{} (Lewis \& Short).
\end{cousinweb}

\begin{grammarlens}[title={What you've built}]
One more word for telling, explaining and describing at B1.
\end{grammarlens}

\subsection*{Your turn}
\begin{itemize}
\raggedright
  \item \emph{Say it:} \emph{contemporaneamente}
  \item \emph{Say it:} \emph{contemporaneamente}, once more
  \item \emph{Say it:} a sentence of your own with \emph{contemporaneamente}
  \item \emph{Recall:} say the Italian for after which, then the Italian for in the meantime, then say \emph{contemporaneamente} again
\end{itemize}

\begin{tcolorbox}[breakable,colback=teal!4,colframe=teal!35!black,title={Before you move on}]
What does contemporaneamente mean? (\textquotedblleft{}At the same time, simultaneously\textquotedblright{}.) Say it once more.
\end{tcolorbox}
\section[non appena]{non appena — as soon as}

\label{lesson:IT-C274-non-appena}



\begin{quote}
Before the new one: say the Italian for in the meantime, then the Italian for at the same time.
\end{quote}

\subsection*{You'll want to know: non appena}
\textbf{non appena} — \textquotedblleft{}as soon as\textquotedblright{}.

\begin{quote}
Non appena sono arrivato, ha cominciato a piovere.
\end{quote}

\begin{quote}
\emph{As soon as I arrived, it started to rain.}
\end{quote}

It means the same as \emph{appena} here; the \emph{non} adds nothing negative.

\begin{cousinweb}
\emph{appena} is inherited from Latin \emph{ad paene}, built on \emph{paene} \textquotedblleft{}almost\textquotedblright{} (EtymDB); \emph{non} comes from Latin \emph{nōn} \textquotedblleft{}not\textquotedblright{}.
\end{cousinweb}

\begin{grammarlens}[title={What you've built}]
One more word for telling, explaining and describing at B1.
\end{grammarlens}

\subsection*{Your turn}
\begin{itemize}
\raggedright
  \item \emph{Say it:} \emph{non appena}
  \item \emph{Say it:} \emph{non appena}, once more
  \item \emph{Say it:} a sentence of your own with \emph{non appena}
  \item \emph{Recall:} say the Italian for in the meantime, then the Italian for at the same time, then say \emph{non appena} again
\end{itemize}

\begin{tcolorbox}[breakable,colback=teal!4,colframe=teal!35!black,title={Before you move on}]
What does non appena mean? (\textquotedblleft{}As soon as\textquotedblright{}.) Say it once more.
\end{tcolorbox}
\section[a un certo punto]{a un certo punto — at a certain point, at some point}

\label{lesson:IT-C274-a-un-certo-punto}



\begin{quote}
Before the new one: say the Italian for at the same time, then the Italian for as soon as.
\end{quote}

\subsection*{You'll want to know: a un certo punto}
\textbf{a un certo punto} — \textquotedblleft{}at a certain point, at some point\textquotedblright{}.

\begin{quote}
A un certo punto la luce si è spenta.
\end{quote}

\begin{quote}
\emph{At a certain point the light went out.}
\end{quote}

A favourite phrase of storytellers to bring in a turn in the story.

\begin{cousinweb}
\emph{punto} is inherited from Latin \emph{punctum} \textquotedblleft{}a point, a prick\textquotedblright{}, from \emph{pungere} \textquotedblleft{}to prick\textquotedblright{} (EtymDB); \emph{certo} comes from Latin \emph{certus} \textquotedblleft{}fixed, sure\textquotedblright{}.
\end{cousinweb}

\begin{grammarlens}[title={What you've built}]
One more word for telling, explaining and describing at B1.
\end{grammarlens}

\subsection*{Your turn}
\begin{itemize}
\raggedright
  \item \emph{Say it:} \emph{a un certo punto}
  \item \emph{Say it:} \emph{a un certo punto}, once more
  \item \emph{Say it:} a sentence of your own with \emph{a un certo punto}
  \item \emph{Recall:} say the Italian for at the same time, then the Italian for as soon as, then say \emph{a un certo punto} again
\end{itemize}

\begin{tcolorbox}[breakable,colback=teal!4,colframe=teal!35!black,title={Before you move on}]
What does a un certo punto mean? (\textquotedblleft{}At a certain point, at some point\textquotedblright{}.) Say it once more.
\end{tcolorbox}
\section[d'un tratto]{d'un tratto — all of a sudden}

\label{lesson:IT-C274-d-un-tratto}



\begin{quote}
Before the new one: say the Italian for as soon as, then the Italian for at a certain point.
\end{quote}

\subsection*{You'll want to know: d'un tratto}
\textbf{d'un tratto} — \textquotedblleft{}all of a sudden\textquotedblright{}.

\begin{quote}
D'un tratto tutti si sono alzati.
\end{quote}

\begin{quote}
\emph{All of a sudden everybody stood up.}
\end{quote}

Also \emph{tutto d'un tratto}; storytellers love it.

\begin{cousinweb}
\emph{tratto} is inherited from Latin \emph{tractus} \textquotedblleft{}a pulling, a stretch\textquotedblright{}, from \emph{trahere} \textquotedblleft{}to draw, to drag\textquotedblright{} (EtymDB): in one pull, in one go.
\end{cousinweb}

\begin{grammarlens}[title={What you've built}]
One more word for telling, explaining and describing at B1.
\end{grammarlens}

\subsection*{Your turn}
\begin{itemize}
\raggedright
  \item \emph{Say it:} \emph{d'un tratto}
  \item \emph{Say it:} \emph{d'un tratto}, once more
  \item \emph{Say it:} a sentence of your own with \emph{d'un tratto}
  \item \emph{Recall:} say the Italian for as soon as, then the Italian for at a certain point, then say \emph{d'un tratto} again
\end{itemize}

\begin{tcolorbox}[breakable,colback=teal!4,colframe=teal!35!black,title={Before you move on}]
What does d'un tratto mean? (\textquotedblleft{}All of a sudden\textquotedblright{}.) Say it once more.
\end{tcolorbox}
